\chapter[Band 25: Zusammenhängende Graphen]{Band 25 auf einen Blick}
\noindent\textbf{Zusammenhängende Graphen}\par\medskip
Vorausgesetzt sind ungerichtete Graphen und die gerichtete Erreichbarkeit
mittels endlicher Walks. Weder Schleifenfreiheit noch ein nichtleerer Träger
werden zusätzlich in die Definition aufgenommen.

\begin{BandUebersicht}
\textbf{Zusammenhang}
\UebFormel{\begin{aligned}
&\ConnectedGraphStruct{V,E}\ \leftrightarrow\\
&\quad\UndirectedGraphStruct{V,E}\land\\
&\quad\forall u,v\in V\quad u\GraphReach E v.
\end{aligned}}
\UebQuellen{\FormulaRefAuto{ConnectedGraphDef}[def]}
&
\textbf{Zugriff auf die Grundstruktur}
\UebFormel{\begin{aligned}
&\ConnectedGraphStruct{V,E}\ \vdash\\
&\quad\UndirectedGraphStruct{V,E}\ \land\\
&\quad\forall u,v\in V\quad u\GraphReach E v.
\end{aligned}}
\UebQuellen{\FormulaRefAuto{ConnectedGraphAccess}[thm]} \\
\addlinespace[.8ex]

\textbf{Verbindung konkreter Endpunkte}
Erreichbarkeit erlaubt gleiche Endpunkte ausdrücklich:\UebFormel{\begin{aligned}
&u,v\in V,\quad u=v\ \lor\\
&\exists n\in\mathbb N\,\exists p\
 (\DirectedWalk V E p n\\
&\hspace{6em}\land p(0)=u\land p(n)=v).
\end{aligned}}
\UebQuellen{\FormulaRefAuto{DirectedReachabilityDef}[def]}
&
\textbf{Beliebige Knoten sind erreichbar}
\UebFormel{\begin{aligned}
&\ConnectedGraphStruct{V,E},\
 u,v\in V\ \vdash\\
&\qquad u\GraphReach E v.
\end{aligned}}
\UebQuellen{\FormulaRefAuto{ConnectedGraphVertexReachability}[thm]} \\
\addlinespace[.8ex]
\end{BandUebersicht}

